\section{A new theorem on a familiar foundation}\label{ch17:sec:start}
Imagine meeting the following sentence in a research paper:
\begin{quote}
Every sequence in a complete, totally bounded metric space has a convergent subsequence.
\end{quote}
Most of its words are familiar: Chapter~\ref{ch:limits} defined metric spaces, convergence, Cauchy sequences and completeness, and Chapter~\ref{ch:contraction} used them all. Two expressions are new: ``totally bounded'' and ``subsequence.''

The conclusion is a form of what mathematicians call \emph{compactness}, which topology studies in great generality. Instead of first taking a course on that theory, ask what the two new ideas do in this particular argument, and learn exactly that much. A first exchange with an AI assistant might go like this.

\textbf{Reader:} I know metric spaces, Cauchy sequences and completeness. Describe how the proof of this statement works, without quoting any theorem I may not know.

\textbf{Assistant:} Keep infinitely many terms of the sequence inside a small region, then inside a smaller one, and so on. Pick later and later terms from these nested selections. The picked terms form a Cauchy sequence, and completeness supplies their limit.

\textbf{Reader:} Which facts do those steps rely on that I may not have met yet?

\textbf{Assistant:} Three. You need to know that the space can be covered by finitely many small balls; that if infinitely many indices are split into finitely many groups, some group still contains infinitely many of them; and how to choose indices that keep increasing.

The reader's side of this exchange is a \emph{minimal-background request}: it says what the reader already knows and asks only for what one particular proof needs. The words ``those steps'' pin down which proof, and that matters: a different proof of the same theorem could need different concepts. A good request also asks to see, right after each new concept, the step of the proof that uses it, which keeps the conversation tied to one argument instead of drifting into a survey of a subject.

You need not follow every phrase of the assistant's compressed answers yet. Together they propose a \emph{dependency map} of the proof (Section~\ref{ch12:sec:theorem}): the steps of the proof and the facts that each step relies on. Once we mark which entries we already know, we call it a \emph{prerequisite map}; its new entries are all we need to learn. This proposed map is incomplete, as Section~\ref{ch17:sec:map} will show. Sections~\ref{ch17:sec:balls}, \ref{ch17:sec:pigeonhole} and~\ref{ch17:sec:subsequence} supply the three missing pieces, each in a short self-contained lesson that we call a \emph{capsule}; the last also sets out the plan of the proof, which Section~\ref{ch17:sec:proof} carries out.
